\subsection{SpatialCPA-v14} \label{subsec:spatialcpav14}

SpatialCPA-v14 is a reference-based method for virtual section generation, implemented in Python as
the package \texttt{spatialcpav14}. It synthesises a two-dimensional tissue section at an arbitrary
depth $z$ from a stack of real sections, by conditional flow matching in a joint
molecular--morphological latent space. The input for v14 includes: (i) a spatial transcriptomics
volume of cells by genes, with per-cell coordinates $(x, y, z)$ in micrometres and, optionally, cell
type labels, and (ii) one scalar target depth $z$ per section to be synthesised. The method operates
in a generation-only regime: the held-out sections are physically absent from the input file, so the
method never observes the cells it is asked to reconstruct, and receives no information about them
beyond the scalar $z$. Expression is log-normalised on the input side according to the assay: raw
counts are library-normalised to a common total of $10^{4}$ and then $\log(1+x)$-transformed, whereas
fluorescence intensities are $\log(1+x)$-transformed only, and matrices that are already normalised
are passed through unchanged. The sections are ordered by their median $z$ into a stack, and all
subsequent computation uses only sections in that stack.

Two per-cell representations are computed first. The expression latent is obtained by standardising
each gene and then projecting onto the leading $d_{e} = 32$ principal axes by truncated SVD, giving a
compact code $\textrm{\textbf{e}}$ that is smooth enough to model with a flow and approximately
invertible. The morphology channels are a per-cell pseudo-image computed \textit{within} a single
section: for each cell, its $k = 12$ nearest in-plane neighbours are weighted by
$w = \exp\left(-\left(d / (\sigma \, \tilde{d})\right)^{2}\right)$, where $\tilde{d}$ is the median
neighbour distance and $\sigma = 1$, yielding a soft local cell-type composition (one column per type)
together with a local density channel defined as the inverse mean neighbour distance, robustly scaled
by its own median. The two representations are fused by an encoder into a joint latent
$\textrm{\textbf{h}}$ of dimension $d = 48$ (hidden width 128, dropout 0.05), which also carries
decoders back to $\textrm{\textbf{e}}$, to the morphology channels, and to a cell-type head.

Conditioning is provided by a three-dimensional positional-attention module: for each query position,
tokens are drawn from the $n_{\textrm{context}} = 16$ nearest real cells plus one global summary token
per section, over the nearest real section on each side of $z$ (\texttt{context\_slices\_each\_side}
$= 1$) together with both flanking sections, with $(x, y, z)$ encoded by 6 Fourier bands and attended
by 4 heads at model width 96. The generative component is a velocity field (hidden width 192, 4
layers, time embedding 32) trained by conditional flow matching on the optimal-transport
straight-line path. Training proceeds in two phases: Phase A fits the encoder and its decoders by
reconstruction for 60 epochs, and Phase B fits the velocity field and the attention module for 160
epochs, using AdamW at learning rate $3\times 10^{-4}$, weight decay $10^{-4}$, gradient clipping at
2.0, and batches of 256 cells. Two perturbations make the model gap-aware: a whole context section is
dropped with probability 0.35, and the query depth is jittered by a Gaussian of scale 0.15 (in units
of the $z$ half-range) so that the model marginalises over depth uncertainty. Phase B additionally
carries biology-informed regularisers, annealed in over the first 40 epochs: a closed-loop
consistency term at weight 0.10 that penalises the discrepancy between a latent and its
re-encoding after decoding; a latent total-variation term, relaxed at morphological edges, whose
smoothness and interface weights (0.05 and 0.10) enter as a single summed coefficient; and a
soft-margin hypoxia term at weight 0.05 with margin 0.02, which enforces the directionality of the
core-to-periphery gradient.

Generation begins by selecting the two real sections that bracket the requested depth. Writing
$z_{lo}$ and $z_{hi}$ for their centres, the interpolation fraction is
$t = (z - z_{lo}) / (z_{hi} - z_{lo})$ and the number of cells to synthesise is
$n = \textrm{round}\left((1-t) \, n_{lo} + t \, n_{hi}\right)$. Positions are then drawn by resampling
the two flanking sections' real coordinates rather than being generated. To prevent the two flanks
from being interleaved cell by cell, which would scramble local neighbourhoods, v14 builds a smooth
random field over the plane as the sum of three random sinusoids with frequencies drawn uniformly on
$[0.5, f]$ and random phases, where $f = 4.0$, rescaled to $[0, 1]$. The lower section supplies the
region where the field is at least $t$ and the upper section supplies the remainder, so the expected
area taken from the upper section is $t$ and composition is still interpolated correctly, while each
generated cell's entire neighbourhood is drawn from a single real section. Positions are sampled
without replacement whenever the pool is large enough, and perturbed by Gaussian jitter of scale
$0.05 \, \tilde{d}$. The procedure returns both the position and the index of the real cell it was
drawn from, which serves as that cell's inherited donor.

Each generated position is then passed through the attention module and the flow: the ODE is
integrated from noise by 12 Euler steps, repeated from 4 independent initial noises whose endpoints
are averaged, and the resulting joint latent is decoded to $\hat{\textrm{\textbf{e}}}$, a code in the
same 32-dimensional space as the real cells. Expression is not generated from this code. Instead, each
cell is grounded in a real profile: with \texttt{ground\_blend\_flow} set to 1.0 every cell is
reconsidered, and for each one the \texttt{ground\_k} spatially nearest real cells of the two flanking
sections form its candidate set, from which the method selects
$\arg\min_{c} \lVert \textrm{\textbf{e}}_{c} - \hat{\textrm{\textbf{e}}} \rVert$ under the default
\texttt{selection="flow"} rule. The candidate set is therefore spatial and the latent only ranks
within it; with \texttt{ground\_k} $= 2$ the flow contributes one binary choice per cell, and because
v14 has no swap-margin parameter, the inherited donor is always overridden. The cell type is inherited
from the selected donor (\texttt{type\_placement="exemplar"}). Composition matching then nudges the
synthesised type mix toward the interpolated flanking composition $(1-t) \, \textrm{comp}(lo) + t \,
\textrm{comp}(hi)$, by re-grounding cells currently in over-represented types to a nearby real cell of
the needed type, again chosen by minimum latent distance among the same \texttt{ground\_k} candidates.

Finally, v14 optionally blends the flow-decoded profile into the emitted expression with weight
\texttt{edit\_weight}; when this weight is set to 0 the blend is skipped and every emitted profile is
one real cell's measurement, copied in full. The result is clipped below at zero and, because the
evaluator expects count-like values, exponentiated as $\textrm{expm1}\left(\textrm{clip}(\cdot, 0,
20)\right)$. It should be noted that this exponentiation is not the inverse of the input
normalisation on count data, because the per-cell library rescaling applied before $\log(1+x)$ is not
undone; on assays where the input transform is $\log(1+x)$ alone, such as fluorescence intensities,
the round trip is exact. v14 returns a virtual section consisting of: (i) the cell coordinates
$(x, y, z)$, (ii) the cell-by-gene expression matrix, and (iii) the cell type label and index per
cell.


\subsection{SpatialCPA-v18} \label{subsec:spatialcpav18}

SpatialCPA-v18 extends v14 with four modifications, each of which targets a specific deficiency
observed in benchmarking, and is implemented as a single-file Python module. Its architecture,
hyperparameters and training procedure are identical to those of v14 described in
Section~\ref{subsec:spatialcpav14}: the same expression latent and morphology channels, the same joint
encoder, the same three-dimensional attention and velocity field, and the same two-phase training
schedule with the same regularisers. All four modifications act at generation time. The input is the
same as v14's with one addition: v18 retains the pre-normalisation expression matrix alongside the
log-normalised one, so that a copied profile can be emitted on its original measurement scale. The
selection of flanking sections, the interpolation fraction $t$, the target cell count, and the
coherent-patch layout are unchanged from v14.

The first modification concerns how a donor is selected. Rather than taking the latent-nearest
candidate unconditionally, v18 protects the inherited donor with a margin: a swap is performed only if
$d_{\textrm{inherited}} - \min_{c} d_{c} \geq \tau$, where $d$ denotes distance in the latent space and
$\tau$ is \texttt{ground\_keep\_margin}, so that a cell whose inherited profile is already close to the
generated latent is left alone. When a swap does occur, the donor is not taken by $\arg\min$ but
sampled, with probability proportional to the product of three terms,
\[
p_{c} \;\propto\; \exp\!\left(-\frac{d_{c} - \min_{c'} d_{c'}}{T}\right)
\cdot \exp\!\left(-\frac{1}{2}\left(\frac{\delta_{c}}{2 \, \bar{s}}\right)^{2}\right)
\cdot \frac{1}{1 + \lambda \, u_{c}},
\]
in which $T$ is the temperature \texttt{ground\_temp}, $\delta_{c}$ is the in-plane distance from the
generated position to candidate $c$ and $\bar{s}$ the median nearest-neighbour spacing of the donor
pool, and $u_{c}$ is the number of times candidate $c$ has already been used, with $\lambda = 0.5$.
The three terms are, respectively, latent similarity, a spatial-locality prior that keeps the copy
local as the candidate set widens, and a reuse penalty. The reuse penalty addresses a specific
failure of the $\arg\min$ rule at small candidate counts, in which many generated cells collapse onto
the same few source cells and the emitted distribution carries high-multiplicity duplicate profiles,
which is precisely what an optimal-transport comparison against the held-out cells penalises. The
usage counter is initialised from the inherited donors and shared by every subsequent re-grounding
pass.

The second modification replaces inherited cell types with a vote. Under
\texttt{type\_mode="vote"}, each generated cell's type is assigned by a distance-weighted $k$-nearest
neighbour vote, with weights $1/(d + 10^{-6})$, over the \texttt{type\_vote\_k} nearest real cells
drawn from \textit{both} flanking sections. Cells whose voted type differs from the type of their
inherited donor are re-grounded to a nearby real cell of the voted type, selected among the same
spatial candidates by latent distance under the sampling rule above, so that the emitted profile
remains consistent with the emitted label. The motivation is that pure inheritance copies each cell's
type from a single source cell, so that along the seam between two coherent patches the local type
mosaic is discontinuous, which the neighbourhood-enrichment metrics read as noise. Composition
matching against the interpolated flanking composition is then applied as in v14.

The third modification introduces partial gene resampling. For each generated cell, a second donor is
drawn from the $\max(\texttt{ground\_k}, 4)$ nearest real cells that share the cell's voted type and
are not its own exemplar; the partner is selected uniformly at random among those candidates, and
cells for which no valid partner exists are left unmixed. A Bernoulli mask is then drawn independently
for each (cell, gene) pair with success probability \texttt{gene\_mix\_frac}, and the masked entries
take their value from the partner rather than the exemplar. The rationale is that whole-profile copies
produce an expression cloud consisting of training atoms, whereas the held-out cells are novel points
on the same manifold, so an optimal-transport comparison systematically favours a method that emits
novel but realistic combinations; mixing a minority of genes from a same-type, same-niche neighbour
creates such combinations while bounding the resulting chimerism. Because both donors are real cells
of the same type in the same neighbourhood, every emitted value remains a real measurement.

The fourth modification concerns the output scale. When the decode blend is disabled and the raw
matrices are available, v18 discards the log-scale array entirely and emits the selected donor's
\textit{raw} measurement, re-applying the gene-mix mask on the raw scale so that mixed entries take
the partner's raw value. The exponentiation used by v14 is therefore never executed. This addresses
the observation that $\textrm{expm1}$ is not the inverse of the input normalisation on count data, and
that on assays with a wide dynamic range the resulting non-linear round trip inverts the ranking of
genes by variance. Like v14, v18 returns a virtual section consisting of: (i) the cell coordinates
$(x, y, z)$, (ii) the cell-by-gene expression matrix, and (iii) the cell type label and index per
cell. The distinction between the two methods at the level of a single emitted cell is that v14
returns one real cell's profile across all genes, transformed, whereas v18 returns a per-gene chimera
of two real cells of the same type, untransformed.
